% from main.tex
%%%%%%%%%%  scan idx 385  |  folio 372  %%%%%%%%%%
résolution partielle globale~:
\[
  L_i\longrightarrow L_{i-1}\longrightarrow\cdots\longrightarrow L_0\longrightarrow M\longrightarrow 0,
\]
par des \(B\)-Modules localement libres de type fini, le noyau \(Z_i\) de \(L_i\to L_{i-1}\) est encore un \(B\)-Module de type fini (I. 2.11). Il est même de présentation finie (en faisant le même raisonnement avec \(i+1\) au lieu de \(i\)). Enfin, il est plat sur \(S\). Il vérifie donc les mêmes hypothèses que \(M\), ce qui permet de construire de proche en proche une résolution globale graduée de \(M\) par des \(B\)-Modules gradués localement libres de type fini, d'après la remarque initiale. Pour \(i\geq r+1\), \(Z_i\) est un \(B\)-Module projectif sur tout ouvert affine (1.5.2) et par suite est localement libre en tant que \(B\)-Module gradué, d'après 1.5.1. On peut alors, dans la résolution graduée construite, remplacer \(L_{r+1}\) par \(Z_{r+1}\), et les termes suivants par \(0\), d'où la résolution annoncée.

Montrons enfin le point c) de la démonstration.

\textbf{Proposition 1.7.} Soient \(S\) un schéma quasi-compact, \(E\) un \(\mathcal O_S\)-Module localement libre de type fini, \(B=\Sym(E)\), et \(L\) un \(B\)-Module gradué, localement libre de type fini (1.2). Alors il existe une filtration décroissante de \(L\) par des sous-\(B\)-Modules gradués \(L^{(m)}\), telle que :
\[
\begin{array}{ll}
\text{i)}& L^{(m)}=L\quad\text{pour }m\leq m_0;\\[1mm]
\text{ii)}& L^{(m)}=0\quad\text{pour }m\geq m_1;\\[1mm]
\text{iii)}& \text{pour tout }m,\quad \operatorname{gr}^m(L)=L^{(m)}/L^{(m+1)}\\
&\quad \text{est isomorphe comme }B\text{-Module gradué à un}\\
&\quad B\text{-Module gradué de la forme }N_m\otimes_{\mathcal O_S}B(m),\\
&\quad \text{où }N_m\text{ est un }\mathcal O_S\text{-Module localement libre de type fini.}
\end{array}
\]

Comme \(S\) est quasi-compact, on peut trouver des entiers \(n_0\) et \(n_1\) tels qu'au voisinage de chaque point de \(S\), les éléments d'une base homogène de \(L\)
